% Lösungen Einheit 5 von 5 – Dreiecke und Figuren aus Tangente, Normale und Achsen (zusammenbau v0.1)
\einheitenkopf{Einheit 5 von 5 · Dreiecke und Figuren aus Tangente, Normale und Achsen}

\erg{23}{a) Umfang – alle drei Seiten, die Hypotenuse mit dem Satz des Pythagoras. \quad b) Nullstelle $x = 2$, $y$-Achsenabschnitt $g(0) = 6$; $A = \frac{1}{2} \cdot 2 \cdot 6 = 6$ \quad c) $h(-1) = 4$, $h'(x) = \frac{4}{x^2}$, $h'(-1) = 4$; $t(x) = 4x + 8$; Achsenabschnitte $x = -2$ und $t(0) = 8$; $A = \frac{1}{2} \cdot 2 \cdot 8 = 8$ (mit Beträgen, nicht $-8$) \quad d) $f(0) = 3$, $f'(0) = -2$, $t(x) = -2x + 3$; Ecken $(0 | 0)$, $(1{,}5 | 0)$, $(0 | 3)$; Hypotenuse $\sqrt{1{,}5^2 + 3^2} \approx 3{,}35$; $U \approx 1{,}5 + 3 + 3{,}35 = 7{,}85$; gleicher Abstand: Mitte der Hypotenuse $(0{,}75 | 1{,}5)$ (Thaleskreis, nicht der Schwerpunkt) \quad e) $f'(x) = \frac{2x}{x^2 + 3}$, $f'(1) = 0{,}5$; $y$-Achsenabschnitte: Tangente $\mathrm{ln}\,4 - 0{,}5$, Normale $\mathrm{ln}\,4 + 2$; Grundseite auf der $y$-Achse $2{,}5$, Höhe $1$ (Abstand von $B$ zur $y$-Achse); $A = \frac{1}{2} \cdot 2{,}5 \cdot 1 = 1{,}25$ \quad f) $f(\pm 2) = 3$, $f'(2) = -1$, $f'(-2) = 1$; $t_1(x) = -x + 5$, $t_2(x) = x + 5$; Ecken $(-5 | 0)$, $(5 | 0)$, $(0 | 5)$; Seiten $10$ und zweimal $\sqrt{50} \approx 7{,}07$; $U \approx 24{,}14$ \quad g) $f'(x) = e^{\frac{1}{2}x}$, $f'(0) = 1$; $t(x) = x + 3$; Achsenschnittpunkte $(-3 | 0)$ und $(0 | 3)$: beide Katheten sind $3$ lang, der rechte Winkel liegt im Ursprung – gleichschenklig (berechnet, nicht abgelesen). \quad h) $h'(0) = -2$, $t(x) = -2x$; Dreieck unterhalb der $x$-Achse: $A = \frac{1}{2} \cdot b \cdot |t(b)| = b^2 = 49$, also $b = 7$ (Fläche positiv ansetzen) \quad i) $f_a'(x) = 2(x - a)^2 + 4x(x - a)$, $f_a'(0) = 2a^2$; $h(x) = -\frac{1}{2a^2}x$ (die Normale, nicht die Tangente); Breite $a$ (Nullstellen $0$ und $a$), Höhe $|h(a)| = \frac{1}{2a}$; $A = a \cdot \frac{1}{2a} = \frac{1}{2}$ – unabhängig von $a$.}
\erg{24}{a) Tom setzt den negativen Achsenabschnitt ohne Betrag ein. Richtig: $A = \frac{1}{2} \cdot 2 \cdot 6 = 6$. \quad b) Zwei Seiten liegen auf den Achsen, und die Achsen stehen im Ursprung senkrecht aufeinander.}
